% ATTEMPT_STATUS: unresolved
% TOP_PROBLEM: {"problemId": "problem.selberg-1-4-conjecture", "canonicalTitle": "Selberg 1/4 Conjecture", "releaseRank": 101, "primaryDomain": "number_theory_arithmetic_geometry", "primaryDomainLabel": "Number theory & arithmetic geometry", "displayStatus": "open", "releaseStatus": "open"}
% !TeX program = lualatex
% Complete problem section copied verbatim from research_batch_101_103.tex.
\documentclass[11pt]{article}
\usepackage[margin=25mm]{geometry}
\usepackage{fontspec}
\setmainfont{Latin Modern Roman}
\usepackage{amsmath,amssymb,mathtools}
\usepackage{unicode-math}
\setmathfont{Latin Modern Math}
\usepackage{enumitem,needspace}
\usepackage{xurl,hyperref}
\hypersetup{colorlinks=true,urlcolor=blue}
\setcounter{secnumdepth}{1}
\setlength{\parindent}{0pt}
\setlength{\parskip}{5pt}
\setlength{\emergencystretch}{3em}
\setlist[itemize]{leftmargin=3em,itemsep=5pt,topsep=4pt,parsep=0pt}
\allowdisplaybreaks
\newcommand{\N}{\mathbb{N}}
\newcommand{\Z}{\mathbb{Z}}
\newcommand{\Q}{\mathbb{Q}}
\newcommand{\R}{\mathbb{R}}
\newcommand{\C}{\mathbb{C}}
\newcommand{\F}{\mathbb{F}}
\newcommand{\A}{\mathbb{A}}
\newcommand{\PP}{\mathbb P}
\newcommand{\E}{\mathbb{E}}
\newcommand{\eps}{\varepsilon}
\newcommand{\dd}{\,\mathrm d}
\newcommand{\sref}[2]{\hyperlink{source:#1:#2}{[S#2]}}
\newcommand{\eref}[2]{\hyperlink{extra:#1:#2}{[E#2]}}
% Turn the next switch off for a shorter reading copy. The quotations preserve
% every exactTarget field, including qualifications omitted from a short summary.
\newif\ifshowcatalogue
\showcataloguetrue
\newcommand{\cataloguescope}[1]{\ifshowcatalogue
 \paragraph{Exact scope in the supplied catalogue.}
 \begin{quote}\small #1\end{quote}\fi}
\begin{document}
\setcounter{section}{100}
% ================================================================
% problemId: problem.selberg-1-4-conjecture
% source releaseRank: 101; source releaseStatus: open
% exactTarget SHA-256: 369a061b0f8d65f7956e94416c425a610ae109ecd8cb0ba4744936d4982da2c1
\clearpage
\section{\texorpdfstring{Selberg 1/4 Conjecture}{Selberg 1/4 Conjecture}}\label{101}
\subsection{Definitions and mathematical statement}
The upper half-plane $\mathbb H=\{x+iy:y>0\}$ has hyperbolic metric $y^{-2}(dx^2+dy^2)$ and nonnegative Laplacian $\Delta=-y^2(\partial_x^2+\partial_y^2)$. A congruence subgroup contains a principal subgroup $\Gamma(N)$ for some $N$. On its finite-area quotient, the assertion excludes any nonconstant $L^2$ eigenfunction with
\[
 \Delta f=\lambda f,\qquad0<\lambda<\tfrac14.
\]
This formulation remains meaningful for noncompact quotients, where a continuous spectrum is also present; ``first nonzero eigenvalue'' is understood as the bottom of the nonconstant spectral range when necessary.
\par\smallskip\textit{Formulation and concept references:} \sref{101}{1}.
\cataloguescope{Prove that the first nonzero eigenvalue lambda\_1 of the Laplace-Beltrami operator on every congruence quotient of the upper half-plane is at least 1/4.}
\subsection{Short English statement}
Do all congruence hyperbolic surfaces have a spectral gap of at least one quarter above the constant functions?
% BEGIN RESEARCH ADDITION ATTEMPT-101
\subsection{Research attempt}

\paragraph{Current frontier.}
Write a possible exceptional eigenvalue as $\lambda=1/4-\nu^2$, with
$0<\nu<1/2$. The uniform Kim--Sarnak bound is
$\nu\le 7/64$, or $\lambda\ge 975/4096$; Pascadi's January 2026 revision
explicitly records this bound in Theorem C. His improved large-sieve estimates
control certain weighted contributions of exceptional forms, rather than proving
that exceptional forms do not exist. \eref{101}{1}
Booker, Lee, and Str\"ombergsson prove the conjecture for $\Gamma_1(N)$,
$N\le880$, and $\Gamma(N)$, $N\le226$, in Theorem 1.1 of
\sref{101}{3}. Neither a finite level range nor a small average contribution
excludes one eigenfunction at an arbitrary remaining level.

The analytic framework used below is the cofinite pre-trace formula, including
the Eisenstein integral, not the trace of an unregularized heat operator on a
noncompact surface. The transform convention and local spectral estimates are
those recalled in Section 2 of \eref{101}{2}. The explicit filter and the
conditional argument that follows are deductions of this notebook; no novelty
claim is made for the general principle of detecting exceptional spectrum by
positive test functions.

\paragraph{Attack route.}
Three routes were compared. Symmetric-power amplification would need suitable
lifts in arbitrarily large degree for the particular Maa\ss{} representation:
if a lift detects $m\nu$ and bounds it by a constant independent of $m$, then
$\nu=0$. This describes a missing functoriality input, not a theorem being
invoked. A level-density argument would need to exclude a single exceptional
representation, not merely show that exceptions have density zero. The route
pursued here instead constructs a spectral filter that kills the constant
function exactly, remains nonnegative on every other spectral component, and
has compact geometric support. Its missing input will be a concrete signed
orbit-counting estimate.

\paragraph{Attempt: reduction to a fixed principal quotient.}
It suffices to handle the effective image of $\Gamma(N)$ in
$\operatorname{PSL}_2(\R)$ for every $N\ge3$. Indeed, any congruence subgroup
contains such a principal subgroup after replacing its level by a sufficiently
large multiple. An exceptional eigenfunction pulls back to the finite cover;
its eigenvalue is unchanged and its squared norm is multiplied by the degree
of the cover. The pullback is still nonconstant and square integrable. Thus
fix a torsion-free principal group $\Gamma$, put
$Y=\Gamma\backslash\mathbb H$, and write $V=\operatorname{vol}(Y)$.
Constants below may depend on this fixed quotient. No uniform-in-level
constant will be silently assumed.

Choose $\phi\in C_c^\infty(Y)$, $\phi\ge0$, positive on a nonempty open set,
and put $\Phi=\int_Y\phi\,d\mu$. Let $u_j$ be an orthonormal list of discrete
$L^2$ eigenfunctions, including the constant and any residual eigenfunctions,
with $\lambda_j=1/4+r_j^2$. The parameter $r_j$ is real and nonnegative,
or lies on $i(0,1/2]$. For $T\ge1$, set
\begin{equation}\label{ra101:filter}
 h_T(r)=\left(r^2+\frac14\right)^2
          \left(\frac{\sin(Tr)}{Tr}\right)^8,
 \qquad \frac{\sin 0}{0}:=1.
\end{equation}
This is entire and even. For real $r$ it is nonnegative. Moreover,
\begin{equation}\label{ra101:exception}
 h_T(i/2)=0,\qquad
 h_T(i\nu)=\left(\frac14-\nu^2\right)^2
             \left(\frac{\sinh(T\nu)}{T\nu}\right)^8>0
       \quad(0<\nu<1/2).
\end{equation}
The zero at $i/2$ removes the constant without subtracting a large positive
quantity numerically. The eighth power provides enough decay on the real
spectral axis after multiplication by the degree-four factor.

\paragraph{A localized positive quantity.}
Let $k_T(s)$ be the radial inverse Selberg transform of $h_T$, where $s$
is hyperbolic distance. Define
\begin{equation}\label{ra101:quantity}
 Q_{\Gamma,\phi}(T)=
 \int_Y\phi(z)\sum_{\gamma\in\Gamma}
                k_T\bigl(d(z,\gamma z)\bigr)\,d\mu(z).
\end{equation}
The cofinite pre-trace formula, specialized to this filter, gives
\begin{align}\label{ra101:pretrace}
 Q_{\Gamma,\phi}(T)
 &=\sum_j h_T(r_j)\int_Y\phi|u_j|^2\,d\mu \\
 &\quad+\frac1{4\pi}\sum_{\mathfrak a}
   \int_{\R}h_T(r)
      \left(\int_Y\phi(z)
        |E_{\mathfrak a}(z,1/2+ir)|^2\,d\mu(z)\right)dr.\notag
\end{align}
Here the sum over $\mathfrak a$ runs over the cusps, with the usual Eisenstein
normalization. This use of the spectral formula and its continuous term is
supported by \eref{101}{2}, Section 2, especially equation (2.2).

For completeness, the needed convergence is not a formal interchange.
The filter is $O_T((1+|r|)^{-4})$ in a fixed horizontal strip. On the compact
support of $\phi$, the cumulative local spectral mass, combining the discrete
and continuous real spectrum up to $|r|\le R$, is $O_{\Gamma,\phi}(1+R^2)$.
Dyadic summation therefore makes the real-spectrum terms absolutely
integrable. There are finitely many exceptional eigenvalues. On the geometric
side $k_T$ has support in $[0,8T]$, as proved below, so only finitely many
orbit terms occur locally, uniformly on this compact support. Thus both sides
of \eqref{ra101:pretrace} are legitimate finite or absolutely convergent
expressions.

All terms on the right of \eqref{ra101:pretrace} are nonnegative. An exceptional
$u$ has
$b_u:=\int_Y\phi|u|^2\,d\mu>0$: otherwise it vanishes on an open set, and
unique continuation for the Laplace equation forces $u=0$. Consequently an
exceptional parameter $\nu$ implies, for all sufficiently large $T$,
\begin{equation}\label{ra101:growth}
 Q_{\Gamma,\phi}(T)\ge
       c_{\Gamma,\phi,u,\nu}\,T^{-8}e^{8\nu T}.
\end{equation}
Neither the constant nor the onset of this estimate is asserted to be uniform
as $\nu\downarrow0$.

Conversely, if no exceptional eigenvalue exists, then $Q_{\Gamma,\phi}(T)$
is bounded independently of $T\ge1$. Indeed, for real $r$,
\[
 0\le h_T(r)\le
 \left(r^2+\frac14\right)^2\min\{1,|r|^{-8}\},
\]
where the right side is interpreted by its bounded extension near zero.
The same local spectral-mass bound integrates this majorant. In particular,
the threshold $r=0$ causes no difficulty: $h_T(0)=1/16$.

\textbf{Filter lemma.} For this fixed $\Gamma$ and $\phi$, absence of
exceptional $L^2$ eigenvalues is equivalent to
\begin{equation}\label{ra101:subexp}
 \forall\epsilon>0\quad
 Q_{\Gamma,\phi}(T)\ll_{\Gamma,\phi,\epsilon}e^{\epsilon T}
 \quad(T\ge1).
\end{equation}
The forward implication is the boundedness just proved. For the reverse
implication, choose $0<\epsilon<8\nu$ in \eqref{ra101:growth}.
\hfill$\square$

This is a proved reformulation, not an improvement of the spectral gap.
Its usefulness depends entirely on estimating \eqref{ra101:quantity} by
arithmetic geometry without first assuming the desired spectral conclusion.

\paragraph{An explicit geometric kernel and a sufficient missing estimate.}
The Fourier convention is $h(r)=\int_{\R}g(u)e^{iru}\,du$. Put
\[
 b_T(u)=\frac{\mathbf 1_{[-T,T]}(u)}{2T},\qquad
 B_T=b_T^{*8},\qquad
 g_T=\left(-\frac{d^2}{du^2}+\frac14\right)^2B_T.
\]
Then $\widehat{B_T}(r)=(\sin(Tr)/(Tr))^8$, so $g_T$ transforms to $h_T$.
The function $B_T$ is a compactly supported piecewise polynomial of class
$C^6$, and $g_T$ is of class $C^2$, supported in $[-8T,8T]$.
Abel inversion of the transform convention in \eref{101}{2}, Section 2,
gives, now expressed in the distance variable,
\begin{equation}\label{ra101:inverse}
 k_T(s)=-\frac1{\sqrt2\,\pi}
       \int_s^\infty\frac{g_T'(u)}{\sqrt{\cosh u-\cosh s}}\,du.
\end{equation}
The square-root singularity is integrable for $s>0$; at $s=0$, evenness gives
$g_T'(u)=O_T(u)$. This formula also proves the stated support property.

Here are bounds adequate for testing the route. Scaling gives
$B_T(u)=T^{-1}B_1(u/T)$, hence
$\|g_T'\|_\infty\ll T^{-2}$ and
$\|g_T''\|_\infty\ll T^{-3}$ for $T\ge1$.
For $s\ge1$, substitute $u=s+v$ in \eqref{ra101:inverse} and use
\[
 \cosh(s+v)-\cosh s
   =2\sinh(v/2)\sinh(s+v/2).
\]
The resulting denominator has an integrable $v^{-1/2}$ singularity at zero
and exponential decay after inversion at infinity. Its logarithmic
$s$-derivative is bounded for $s\ge1$. Differentiation under this dominated
integral is therefore allowed and yields
\begin{equation}\label{ra101:kernelbound}
 |k_T(s)|+|k_T'(s)|\ll T^{-2}e^{-s/2}\quad(s\ge1).
\end{equation}
For $0\le s\le1$, use
$|g_T'(u)|\le\|g_T''\|_\infty u$ for $u\le2$, and the exponential
bound for the remaining integral; this gives $\sup_{[0,1]}|k_T|\ll1$.
No derivative bound at the origin is needed below.

Consider the locally averaged orbit count
\[
 N_\phi(R)=\int_Y\phi(z)\,
          \#\{\gamma\in\Gamma:d(z,\gamma z)\le R\}\,d\mu(z)
\]
and the discrepancy with \emph{only the constant eigenfunction's main term}
removed:
\begin{equation}\label{ra101:error}
 \mathcal E_\phi(R)=N_\phi(R)
       -\frac{\Phi}{V}\,2\pi(\cosh R-1).
\end{equation}
Since
$2\pi\int_0^\infty k_T(R)\sinh R\,dR=h_T(i/2)=0$, one has the exact
Stieltjes identity
\begin{equation}\label{ra101:stieltjes}
 Q_{\Gamma,\phi}(T)=\int_{[0,8T]}k_T(R)\,d\mathcal E_\phi(R).
\end{equation}
The atom at $R=0$, coming from the identity element, is included.

A sufficient arithmetic input would be the following statement:
\begin{equation}\label{ra101:missing}
 \boxed{\quad
 \forall\epsilon>0\qquad
 |\mathcal E_\phi(R)|\ll_{\Gamma,\phi,\epsilon}
           e^{(1/2+\epsilon)R}\quad(R\ge1).
 \quad}
\end{equation}
To check the implication carefully, the part of
\eqref{ra101:stieltjes} on $[0,1]$ is $O_{\Gamma,\phi}(1)$ because the
counting measure has finite mass there and $k_T$ is uniformly bounded.
On $[1,8T]$, integration by parts and \eqref{ra101:kernelbound} give
\[
 |Q_{\Gamma,\phi}(T)|
 \ll_{\Gamma,\phi,\epsilon}
 1+T^{-2}\int_1^{8T}e^{\epsilon R}\,dR
 \ll_{\Gamma,\phi,\epsilon}1+e^{8\epsilon T}.
\]
The outer endpoint term vanishes because $k_T(8T)=0$; the inner endpoint is
bounded. Choosing $\epsilon$ arbitrarily small proves
\eqref{ra101:subexp}. Thus \eqref{ra101:missing}, for every principal quotient,
would imply exactly the conjecture in this entry, including possible residual
spectrum and noncompact quotients.

\paragraph{Stress test.}
\emph{The indispensable sign cancellation.}
It is tempting to discard signs in the orbit sum. That destroys the argument.
A packing bound on compact subsets gives $N_\phi(R)\ll_{\Gamma,\phi}e^R$.
Combining this with $|k_T(R)|\ll T^{-2}e^{-R/2}$ gives at best an exponential
bound of order $e^{4T}$, up to polynomial factors. Comparison with
\eqref{ra101:growth} recovers only the trivial restriction $\nu\le1/2$.
Furthermore, $k_T$ cannot be nonnegative everywhere: its hyperbolic integral
is zero, while its transform is not identically zero. Spectral positivity
and geometric positivity are different assertions.

\emph{A circular subtraction to avoid.}
In the usual hyperbolic circle problem the main term often includes explicit
contributions from all small eigenvalues. This is precisely the convention
in \eref{101}{2}, equation (1.2). A square-root-size error \emph{after}
subtracting those terms would not prove Selberg. Our
\eqref{ra101:error} deliberately does not subtract them. Estimate
\eqref{ra101:missing} is therefore stronger than a remainder estimate that
already allows exceptional spectrum; it has not been proved here.

\emph{Why an average or a limiting argument does not close the gap.}
An exceptional eigenfunction on one quotient persists on every finite cover.
Dividing a spectral count by the growing covering degree can nevertheless
make this fixed contribution invisible. Thus a density-zero statement in a
tower is compatible with the very counterexample being sought. Here $\Gamma$
and $\phi$ are fixed before $T\to\infty$. Likewise, applying the known
$\nu\le7/64$ bound to the positive spectral side only gives a bound with
exponential rate at most $8(7/64)=7/8$; feeding that bound back into
\eqref{ra101:growth} reproduces $7/64$, not a smaller number.
\eref{101}{1}

\emph{Dependency check.}
Improved exceptional-spectrum estimates can feed into the distribution of
prime factors of polynomial values, as Pascadi's application demonstrates.
They do not by themselves distinguish a prime value from a product of two
primes. Thus neither this conditional route nor the existing weighted
estimates establish the Bunyakovsky assertion in Section~\ref{102}.
\eref{101}{1}

\paragraph{Outcome.}
\textbf{Outcome: Conditional reduction.}

The explicit filter, the localized positivity lemma, and the implication
from \eqref{ra101:missing} have been justified. The missing estimate is stated
with its exact main term, quantifiers, and dependence on the quotient. The
attempt also identifies the exponential loss caused by taking absolute values
and the circularity of subtracting unknown exceptional terms.

No new numerical lower bound for $\lambda_1$ has been obtained. A complete
resolution by this route requires an independent arithmetic proof of
\eqref{ra101:subexp}, or of the stronger orbit discrepancy estimate
\eqref{ra101:missing}, on every principal quotient. Positivity alone does not
supply that proof.

% END RESEARCH ADDITION ATTEMPT-101
\subsection{Sources}
\begin{itemize}\raggedright
\item[\textnormal{[S1]}]\hypertarget{source:101:1}{} Formal statement, Chapter 4, Selberg property.
\url{https://www.ma.huji.ac.il/~alexlub/BOOKS/On%20property/On%20property.pdf}.
{\small\textit{Catalogue formulation source.}}
\item[\textnormal{[S2]}]\hypertarget{source:101:2}{} Asymptotic independence for random permutations from surface groups.
\url{https://link.springer.com/article/10.1007/s10711-025-01003-8}.
{\small\textit{Catalogue research/status reference; not a proof endorsement.}}
\item[\textnormal{[S3]}]\hypertarget{source:101:3}{} Twist-minimal trace formulas and the Selberg eigenvalue conjecture.
\url{https://arxiv.org/abs/1803.06016}.
{\small\textit{Catalogue research/status reference; not a proof endorsement.}}
% BEGIN RESEARCH ADDITION REFERENCES-101
\item[\textnormal{[E1]}]\hypertarget{extra:101:1}{}
A. Pascadi, \emph{Large sieve inequalities for exceptional Maass forms and
the greatest prime factor of $n^2+1$}, Forum of Mathematics, Pi \textbf{14}
(2026), e8; arXiv:2404.04239v3, 15 January 2026. Introduction and Theorem C
(the Kim--Sarnak $7/64$ bound), and the exceptional-spectrum large-sieve
results. \url{https://arxiv.org/html/2404.04239v3}.
\url{https://doi.org/10.1017/fmp.2026.10025}.
{\small\textit{Consulted 25 September 2026; weighted estimates are not an
exclusion of exceptional eigenvalues.}}
\item[\textnormal{[E2]}]\hypertarget{extra:101:2}{}
G. Cherubini and M. S. Risager, \emph{On the variance of the error term in
the hyperbolic circle problem}, Revista Matem\'atica Iberoamericana
\textbf{34} (2018), no.~2, 655--685, DOI 10.4171/RMI/1000.
Section 2, transform convention and pre-trace formula (2.2); equation (1.2)
for the main term including small eigenvalues; Section 4 for smoothing and
absolute convergence.
\url{https://ems.press/content/serial-article-files/38698}.
{\small\textit{Consulted 25 September 2026.}}
% END RESEARCH ADDITION REFERENCES-101
\end{itemize}

\end{document}
